\par\Needspace{30\baselineskip}\pdfbookmark[1]{Prasse. 1814 Dezember 19.}{prasse1814}
\begin{center}
\rule{0.63\linewidth}{0.3pt}\\
Göttingische gelehrte Anzeigen. 1814 Dezember 19.\\
\rule{0.63\linewidth}{0.3pt}
\end{center}

\emph{Tables logarithmiques pour les nombres, les sinus et les tangentes, disposées
dans un nouvel ordre par M. \textup{\textsc{de Prasse}}, professeur des mathématiques à Berlin} (zu
Leipzig) \emph{corrigées et précédées d’une introduction traduite de l’allemand et accompagnée de notes et d’un avertissement par} M. \textsc{Halma} 1814. Paris. 80 Seiten. Preis
Ein Franc.

Ein neuer Druck der handlichen Tafeln von \textsc{Prasse}, die wir in diesen Blättern [1811 Mai 25] angezeigt haben. Durch einen um drei Jahre längeren Gebrauch haben wir unser damaliges Urteil über die Anordnung der Tafeln, die der französische Herausgeber unverändert beibehalten hat, in jeder Hinsicht bestätigt gefunden. Was die Schönheit von Druck und Papier betrifft, scheint uns die französische Ausgabe der deutschen eher nachzustehen; mehrere Druckfehler der letzteren sind hier jedoch berichtigt. Auch wenn man bei einem Werk dieser Art, das der verstorbene von \textsc{Prasse} gewiss nicht des Gewinns wegen, sondern im Dienst der Wissenschaft auf eigene Kosten unternommen hat, nicht weiter untersuchen will, inwieweit Herr \textsc{Halma} zu einem neuen Druck berechtigt war, kann man sich doch nur wundern, dass er aus Sorge, seinerseits wieder nachgedruckt zu werden, die einzelnen Exemplare mit seinem Namenszug versehen hat und den Nachdruckern mit gerichtlichen Schritten droht.
\SourceLeaf{253}
